% ============================================================
% APPENDICE C — Comandi e Scorciatoie
% ============================================================
\chapter{Comandi e Scorciatoie}

% ---------------------------------------------------------
\section{VS Code}

\begin{center}
\begin{tabularx}{\textwidth}{l X}
\toprule
\textbf{Scorciatoia} & \textbf{Azione} \\
\midrule
\texttt{Ctrl+Shift+P} / \texttt{Cmd+Shift+P} & Command Palette \\
\texttt{Ctrl+Shift+I} / \texttt{Cmd+Shift+I} & Apri Copilot Chat (Agent Mode) \\
\texttt{Ctrl+I} / \texttt{Cmd+I} & Copilot Inline Chat \\
\texttt{Tab} & Accetta suggerimento Copilot \\
\texttt{Esc} & Rifiuta suggerimento Copilot \\
\texttt{Ctrl+`} & Apri/chiudi terminale integrato \\
\texttt{Ctrl+Shift+E} & Explorer (file) \\
\texttt{Ctrl+Shift+G} & Source Control (Git) \\
\texttt{Ctrl+P} & Quick Open (cerca file) \\
\texttt{Ctrl+Shift+F} & Cerca in tutti i file \\
\bottomrule
\end{tabularx}
\end{center}

% ---------------------------------------------------------
\section{Node.js e npm}

\begin{center}
\begin{tabularx}{\textwidth}{l X}
\toprule
\textbf{Comando} & \textbf{Descrizione} \\
\midrule
\texttt{npm init -y} & Crea package.json \\
\texttt{npm install <pacchetto>} & Installa dipendenza \\
\texttt{npm install -D <pacchetto>} & Installa come dev dependency \\
\texttt{npm run dev} & Avvia in modalità sviluppo \\
\texttt{npm run build} & Build di produzione \\
\texttt{npm audit} & Controlla vulnerabilità \\
\texttt{npm audit fix} & Correggi vulnerabilità \\
\texttt{npx <comando>} & Esegui pacchetto senza installare \\
\bottomrule
\end{tabularx}
\end{center}

% ---------------------------------------------------------
\section{Prisma}

\begin{center}
\begin{tabularx}{\textwidth}{l X}
\toprule
\textbf{Comando} & \textbf{Descrizione} \\
\midrule
\texttt{npx prisma init} & Inizializza Prisma \\
\texttt{npx prisma migrate dev -{}-name <n>} & Crea e applica migrazione (dev) \\
\texttt{npx prisma migrate deploy} & Applica migrazioni (produzione) \\
\texttt{npx prisma generate} & Rigenera il client Prisma \\
\texttt{npx prisma studio} & Interfaccia web per il database \\
\texttt{npx prisma db push} & Sincronizza schema senza migrazione \\
\texttt{npx prisma db seed} & Esegui seed del database \\
\bottomrule
\end{tabularx}
\end{center}

% ---------------------------------------------------------
\section{Git}

\begin{center}
\begin{tabularx}{\textwidth}{l X}
\toprule
\textbf{Comando} & \textbf{Descrizione} \\
\midrule
\texttt{git init} & Inizializza repository \\
\texttt{git add .} & Aggiungi tutti i file modificati \\
\texttt{git commit -m "messaggio"} & Crea commit \\
\texttt{git status} & Stato dei file \\
\texttt{git log -{}-oneline} & Cronologia commit compatta \\
\texttt{git branch <nome>} & Crea branch \\
\texttt{git checkout <nome>} & Cambia branch \\
\texttt{git merge <nome>} & Unisci branch \\
\texttt{git push origin main} & Pusha su remote \\
\bottomrule
\end{tabularx}
\end{center}

% ---------------------------------------------------------
\section{Flutter}

\begin{center}
\begin{tabularx}{\textwidth}{l X}
\toprule
\textbf{Comando} & \textbf{Descrizione} \\
\midrule
\texttt{flutter create <nome>} & Crea nuovo progetto \\
\texttt{flutter run} & Avvia app (emulatore/dispositivo) \\
\texttt{flutter run -d chrome} & Avvia app nel browser \\
\texttt{flutter pub get} & Installa dipendenze \\
\texttt{flutter pub add <pacchetto>} & Aggiungi dipendenza \\
\texttt{flutter test} & Esegui test \\
\texttt{flutter analyze} & Analisi statica del codice \\
\texttt{flutter build apk -{}-release} & Build APK Android \\
\texttt{flutter build appbundle} & Build AAB per Play Store \\
\texttt{flutter build ipa} & Build IPA per App Store \\
\texttt{flutter doctor} & Verifica ambiente di sviluppo \\
\texttt{flutter devices} & Lista dispositivi disponibili \\
\texttt{flutter screenshot} & Cattura screenshot dall'emulatore \\
\bottomrule
\end{tabularx}
\end{center}

% ---------------------------------------------------------
\section{Docker (PostgreSQL)}

\begin{center}
\begin{small}
\begin{tabularx}{\textwidth}{l X}
\toprule
\textbf{Comando} & \textbf{Descrizione} \\
\midrule
\texttt{docker run -{}-name pg -e POSTGRES\_PASSWORD=password -p 5432:5432 -d postgres:16} & Avvia PostgreSQL \\
\texttt{docker start pg} & Riavvia container \\
\texttt{docker stop pg} & Ferma container \\
\texttt{docker exec -it pg psql -U postgres} & Apri console PostgreSQL \\
\bottomrule
\end{tabularx}
\end{small}
\end{center}

% ---------------------------------------------------------
\section{Claude Code}

\begin{center}
\begin{tabularx}{\textwidth}{l X}
\toprule
\textbf{Comando} & \textbf{Descrizione} \\
\midrule
\texttt{claude} & Avvia sessione interattiva \\
\texttt{claude "richiesta"} & Esegui richiesta singola \\
\texttt{claude -{}-model sonnet} & Specifica modello \\
\texttt{/clear} & Pulisci contesto della sessione \\
\texttt{/cost} & Mostra costo della sessione \\
\texttt{/help} & Mostra comandi disponibili \\
\bottomrule
\end{tabularx}
\end{center}

% ---------------------------------------------------------
\section{Comandi di Deploy}

\begin{center}
\begin{tabularx}{\textwidth}{l X}
\toprule
\textbf{Piattaforma} & \textbf{Comando/Azione} \\
\midrule
Railway & \texttt{git push origin main} (deploy automatico) \\
Vercel & \texttt{git push origin main} (deploy automatico) \\
Vercel CLI & \texttt{npx vercel -{}-prod} (deploy manuale) \\
Railway CLI & \texttt{railway up} (deploy manuale) \\
\bottomrule
\end{tabularx}
\end{center}
